Zeigen, Sie dass in einem Differentialkörper die sogenannte
\emph{logarithmische Ableitung}
\index{logarithmische Ableitung}%
$Df/f$ die folgenden Regeln erfüllt:
\begin{teilaufgaben}
\item
Für zwei Funktionen $f,g\in F$ gilt
$\displaystyle \frac{D(fg)}{fg}=\frac{Df}{f}+\frac{Dg}{g}$.
\item
Für zwei Funktionen $f,g\in F$ gilt
$\displaystyle \frac{D(f/g)}{f/g}=\frac{Df}{f}-\frac{Dg}{g}$.
\item
Zwei Funktionen unterscheiden sich genau dann um einen konstanten
Faktor, wenn sie die gleiche logarithmische Ableitung haben.
\end{teilaufgaben}

\begin{loesung}
Beide Teilaufgaben ergeben sich durch direkte Rechnung mit den
Rechenregeln für die Derivation $D$:
\begin{teilaufgaben}
\item Mit der Produktregel folgt:
\begin{align*}
\frac{D(fg)}{fg}
&=
\frac{Df\cdot g + f\cdot Dg}{fg}
=
\frac{Df}{f} + \frac{Dg}{g}
\end{align*}
\item Mit der Quotientenregel folgt:
\begin{align}
\frac{D(f/g)}{f/g}
&=
\frac{g}{f}\cdot \frac{Df\cdot g - f\cdot Dg}{g^2}
=
\frac{Df\cdot g - f\cdot Dg}{fg}
=
\frac{Df}{f} - \frac{Dg}{g}.
\label{buch:702:logabl}
\end{align}
\item
Die Aussage ist gleichbedeutend damit, dass $f/g$ konstant ist genau
dann, wenn $f$ und $g$ die gleiche logarithmische Ableitung haben.
Falls $f/g$ konstant ist, ist $D(f/g)=0$ und damit erst recht die
logarithmische Ableitung von $f/g$, die nach
\eqref{buch:702:logabl}
die Differenz der logarithmischen Ableitungen von $f$ und $g$ ist.
Sind umgekehrt die logarithmischen Ableitungen von $f$ und $g$ gleich,
dann verschwindet die logarithmische Ableitung von $f/g$ und somit auch
die Ableitung von $f/g$, $f/g$ ist konstant.
\qedhere
\end{teilaufgaben}
\end{loesung}
